\documentclass[a4paper, 12pt]{article}

\newcommand{\templates}{../../template}
\usepackage[a4paper, margin=2.5cm]{geometry}

\usepackage{enumitem}
\setlist[itemize]{noitemsep}
\setlist[enumerate]{noitemsep}

\let\oldpar\paragraph
\renewcommand{\paragraph}[1]{\oldpar{#1\\}\noindent}

% Avoid dots in the table of contents, it mess with the gulpease calculation
\makeatletter
\renewcommand{\@dotsep}{10000} 
\makeatother
\usepackage{graphicx}
\usepackage{hyperref}
\usepackage{makecell}
\usepackage{fancyhdr}

\newcommand{\settitolo}[1]{\newcommand{\titolo}{#1\\}}
\newcommand{\setprogetto}[1]{\newcommand{\progetto}{#1\\}}
\newcommand{\setcommittenti}[1]{\newcommand{\committenti}{#1\\}}
\newcommand{\setredattori}[1]{\newcommand{\redattori}{#1\\}}
\newcommand{\setrevisori}[1]{\newcommand{\revisori}{#1\\}}
\newcommand{\setresponsabili}[1]{\newcommand{\responsabili}{#1\\}}
\newcommand{\setversione}[1]{
	\ifdefined\versione\renewcommand{\versione}{#1\\}
	\else\newcommand{\versione}{#1\\}\fi
}
\newcommand{\setdestuso}[1]{\newcommand{\uso}{#1\\}}
\newcommand{\setdescrizione}[1]{\newcommand{\descrizione}{#1\\}}

\newcommand{\makefrontpage}{
	\begin{titlepage}
		\begin{center}

		\includegraphics[width=0.4\textwidth]{\templates/4ourSquared_logo}\\

		{\Large 4OURSQUARED}\\[6pt]
		\href{mailto://4oursquared.unipd@gmail.com}{4oursquared.unipd@gmail.com}\\
		
		\ifdefined\progetto
		\vspace{1cm}
		{\Large\progetto}
		{\large\committenti}
		\else\fi
		
		\vspace{1.5cm}
		{\LARGE\titolo}
		
		\vfill
		
		\begin{tabular}{r | l}
		\multicolumn{2}{c}{\textit{Informazioni}}\\
		\hline
		
		\ifdefined\redattori
			\textit{Redattori} &
			\makecell[l]{\redattori}\\
		\else\fi
		\ifdefined\revisori
			\textit{Revisori} &
			\makecell[l]{\revisori}\\
		\else\fi
		\ifdefined\responsabili
			\textit{Responsabili} &
			\makecell[l]{\responsabili}\\
		\else\fi
		
		\ifdefined\versione
			\textit{Versione} & \versione
		\else\fi
		
		\textit{Uso} & \uso
		
		\end{tabular}
		
		\vspace{2cm}
		
		\ifdefined\descrizione
		Descrizione
		\vspace{6pt}
		\hrule
		\descrizione
		\else\fi
		\end{center}
	\end{titlepage}
}
\usepackage{hyperref}
\usepackage{array}
\usepackage{tabularx}
\usepackage{adjustbox}

\newcounter{verscount}
\setcounter{verscount}{0}
\newcommand{\addversione}[5]{
	\ifdefined\setversione
		\setversione{#1}
	\else\fi
	\stepcounter{verscount}
	\expandafter\newcommand%
		\csname ver\theverscount \endcsname{#1&#2&#3&#4&#5}
}

\newcommand{\listversioni}{
	\ifnum\value{verscount}>1
		\csname ver\theverscount \endcsname
		\addtocounter{verscount}{-1}
		\\\hline
		\listversioni
	\else
		\csname ver\theverscount \endcsname\\\hline
	\fi
}

\newcommand{\makeversioni}{
	\begin{center}
		\begin{tabularx}{\textwidth}{|c|c|c|c|X|}
		\hline
		\textbf{Versione} & \textbf{Data} & \textbf{Redattore} & \textbf{Verificatore} & \textbf{Descrizione} \\
		\hline
		\listversioni
		\end{tabularx}
	\end{center}
	\clearpage
}

\settitolo{Motivazione delle Scelte}
\setprogetto{Lumos Minima}
\setcommittenti{Imola Informatica}
\setredattori{Ceccato Francesco \\ Soldà Matteo}
\setdestuso{esterno}
\setdescrizione{
Questo documento riporta i confronti tra i vari linguaggi presi in considerazione per lo sviluppo del progetto e la motivazione della scelta di uno rispetto agli altri.
}

\addversione{0.0.0}{09/05/2023}{Soldà Matteo}{Ceccato Francesco}{Prima stesura.}


\begin{document}
\makeindexdetails
\makefrontpage \makeversioni
\tableofcontents
\newpage
\clearpage
\makecontentsdetails{Motivazione delle Scelte}
\newpage
\section{Introduzione}
Il seguente documento ha come scopo principale quello di confrontare i vari linguaggi, database e framework con lo scopo di fornire un prodotto quanto più efficace, efficiente e aderente alle richieste del proponente.\\
\subsection{Sitografia}
\begin{itemize}
    \item https://www.quora.com/Why-are-there-so-many-backend-systems-written-in-dynamically-typed-languages-while-the-benefit-of-static-typing-is-obvious
    \item https://mobisoftinfotech.com/resources/blog/typescript-vs-javascript/
    \item https://stackoverflow.com/questions/49640121/mixing-javascript-and-typescript-in-node-js
    \item https://www.monocubed.com/blog/top-python-frameworks/
    \item https://www.theserverside.com/tip/The-differences-between-Java-and-TypeScript-devs-must-know
    \item https://medium.com/rewrite-tech/embedding-typescript-in-java-a343576031be
    \item https://blog.geekandjob.com/framework-javascript/
\end{itemize}

\newpage
\section{Confronti}
\subsection{Database}
Come database, non sapendo al momento come saranno strutturati i dati che necessitano di essere salvati in maniera persistente, abbiamo preso in considerazione \href{https://mariadb.org/}{MariaDB} per i database SQL e \href{https://www.mongodb.com/}{MongoDB} per i database NoSQL. \\
Queste due scelte sono state fatte per due motivi principali:
\begin{itemize}
    \item Ricerca di un database computazionalmente leggero e che fosse aderente agli standard
    \item Ricerca di un database la cui community è attiva e la documentazione chiara e puntuale, così da poter risolvere quanto prima i problemi che si potessero presentare data la poca esperienza accumulata con i database nel tempo
\end{itemize}
\subsection{Backend}
Per quanto riguarda il backend, la scelta sarà vincolata dalla preferenza verso un linguaggio staticamente o dinamicamente tipizzato. \\
Qualora si scegliesse la tipizzazione statica, le scelte ricadono su:
\begin{itemize}
    \item \href{https://www.java.com/it/}{Java}
    \item \href{https://learn.microsoft.com/en-us/dotnet/csharp/}{C\#}
    \item \href{https://www.typescriptlang.org/}{TypeScript}
\end{itemize}
Mentre, scegliendo la tipizzazione dinamica, le scelte ricadrebbero su:
\begin{itemize}
    \item \href{https://www.python.org/}{Python}
    \item \href{https://www.javascript.com/}{JavaScript}
\end{itemize}
\subsubsection{C\#}
\textbf{Pro:}
\begin{itemize}
    \item Facile da eseguire su dispositivi Windows (che rappresenta la stragrande maggioranza degli utenti)
    \item Essendo un linguaggio compilato e non interpretato, rende difficile in caso di attacchi la lettura del codice sorgente in chiaro
\end{itemize}
\vspace*{0.5cm}
\textbf{Contro:}
\begin{itemize}
    \item Deve essere ricompilato per ogni minima modifica
    \item Hostabile solo su server Windows
\end{itemize}

\subsubsection{Java}
\textbf{Pro:}
\begin{itemize}
    \item Linguaggio facile da imparare
    \item Linguaggio di alto livello (assenza di puntatori, presenza di garbage collector)
    \item Linguaggio OOP
    \item \textit{Platform Indipendent}
    \item Supporto al \textit{multithreading\textsubscript{G}} 
\end{itemize}
\vspace*{0.5cm}
\textbf{Contro:}
\begin{itemize}
    \item Rispetto ad altri linguaggi può risultare più lento o povero in fatto di performance
    \item Mancanza di feature per il backup dei dati
    \item Rispetto ad altri linguaggi, richiede uno spazio di memoria significativo
    \item Rispetto ad altri linguaggi, usa molto più codice per svolgere una operazione basilare
\end{itemize}
\subsubsection{JavaScript}
\textbf{Pro:}
\begin{itemize}
    \item Semplice da imparare
    \item Interoperabile
    \item Versatile in quanto linguaggio \textit{fullstack\textsubscript{G}}
\end{itemize}
\vspace*{0.5cm}
\textbf{Contro:}
\begin{itemize}
    \item Linguaggio interpretato
    \item Tipizzazione dinamica
    \item Inconsistenza tra browser
\end{itemize}
\subsubsection{Python}
\textbf{Pro:}
\begin{itemize}
    \item Facile da imparare
    \item Flessibile ed estensibile
    \item Presenza di moltissime librerie
    \item Estremamente scalabile
    \item Portabile
    \item Vasta integrazione con l'\textit{IoT\textsubscript{G}}
\end{itemize}
\vspace*{0.5cm}
\textbf{Contro:}
\begin{itemize}
    \item Forte utilizzo di memoria
    \item Tipizzazione dinamica
    \item Multithreading complesso
    \item Garbace colletion che porta ad un potenziale spreco di memoria
\end{itemize}
\subsubsection{Typescript}
\textbf{Pro:}
\begin{itemize}
    \item Linguaggio compilato
    \item Tipizzazione statica
\end{itemize}
\vspace*{0.5cm}
\textbf{Contro:}
\begin{itemize}
    \item Inconsistenza tra browser
    \item Tipizzazione alcune volte complessa
    \item Necessità di ricompilare per ogni minima modifica
\end{itemize}
\subsection{Backend Framework}
\subsubsection{C\#}
\begin{itemize}
    \item \href{https://dotnet.microsoft.com/en-us/apps/aspnet}{ASP.NET}
\end{itemize}
\subsubsection{Java}
\begin{itemize}
    \item \href{https://spring.io/}{Java Spring}
\end{itemize}
\subsubsection{JavaScript}
\begin{itemize}
    \item \href{https://nextjs.org/}{Next}
    \item \href{https://nodejs.org/}{Node}
\end{itemize}
\subsubsection{Python}
\begin{itemize}
    \item \href{https://www.djangoproject.com/}{Django}
    \item \href{https://flask.palletsprojects.com/}{Flask}
\end{itemize}
\subsubsection{Typescript}
\begin{itemize}
    \item \href{https://nestjs.com/}{Nest}
    \item \href{https://feathersjs.com/}{Feathers}
    \item \href{https://loopback.io/}{Loopback}
\end{itemize}
\subsection{Frontend}
Escludendo \textit{HTML5}, \textit{XHTML} e \textit{CSS} che saranno di sicuro usati per la parte di frontend della webapp, i linguaggi tra cui scegliere risultano essere solamente
\begin{itemize}
    \item \href{https://www.typescriptlang.org/}{TypeScript}
    \item \href{https://www.javascript.com/}{JavaScript}
\end{itemize}
\subsubsection{JavaScript}
\subsubsection{TypeScript}

\subsection{Frontend Framework}
\subsubsection{JavaScript}
\begin{itemize}
    \item \href{https://angularjs.org/}{Angular}
    \item \href{https://react.dev/}{React}
    \item \href{https://vuejs.org/}{Vue}
    \item \href{https://backbonejs.org/}{Backbone}
    \item \href{https://preactjs.com/}{Preact}
    \item \href{https://expressjs.com/}{Express}
\end{itemize}
\subsubsection{TypeScript}
\begin{itemize}
    \item \href{https://angularjs.org/}{Angular}
    \item \href{https://emberjs.com/}{Ember}
\end{itemize}
Per i pro e i contro dei framework, sia frontend che backend, si rimanda il lettore a \href{https://blog.geekandjob.com/framework-javascript/}{questo link}.
\newpage
\section{Scelte e Motivazioni}
[WIP]
\end{document} 

%% TO-DO
% Aggiungere descrizione con pro e contro dei linguaggi
% Aggiungere descrizione con pro e contro dei framework
